\documentclass[UTF8,fontset=fandol,11pt]{ctexart}
\usepackage[a4paper,margin=1.75cm]{geometry}
\usepackage{amsmath,amssymb}
\usepackage{enumitem}
\usepackage{longtable}
\usepackage{booktabs}
\usepackage{tcolorbox}
\tcbuselibrary{skins,breakable}
\usepackage{xcolor}
\usepackage{hyperref}
\usepackage{graphicx}
\usepackage{tabularx}
\usepackage{fancyhdr}
\usepackage{titlesec}
\usepackage{array}
\hypersetup{colorlinks=true,linkcolor=blue,urlcolor=blue}
\pagestyle{fancy}
\fancyhf{}
\lhead{\small 工程管理复习讲义}
\rhead{\small \thepage}
\setlength{\parindent}{0pt}
\setlist[enumerate]{leftmargin=2em,itemsep=0.3em}
\setlist[itemize]{leftmargin=1.8em,itemsep=0.2em}
\definecolor{mainblue}{HTML}{2454D6}
\definecolor{deepblue}{HTML}{13204B}
\definecolor{softblue}{HTML}{EEF5FF}
\definecolor{mainpurple}{HTML}{7B2CBF}
\definecolor{softpurple}{HTML}{F5ECFF}
\definecolor{mainorange}{HTML}{F77F00}
\definecolor{softorange}{HTML}{FFF3E0}
\definecolor{maingreen}{HTML}{138A36}
\definecolor{softgreen}{HTML}{EAF8ED}
\definecolor{answerred}{HTML}{D62828}
\newtcolorbox{mainbox}[1][]{enhanced,colback=softblue,colframe=mainblue,boxrule=0.9pt,arc=2mm,left=2.5mm,right=2.5mm,top=1.5mm,bottom=1.5mm,#1}
\newtcolorbox{quizbox}[1][]{enhanced,breakable,colback=white,colframe=mainpurple,coltitle=white,colbacktitle=mainpurple,fonttitle=\bfseries,boxrule=0.8pt,arc=2mm,left=2.8mm,right=2.8mm,top=1.6mm,bottom=1.6mm,#1}
\newtcolorbox{answerbox}[1][]{enhanced,breakable,colback=softgreen,colframe=maingreen,coltitle=white,colbacktitle=maingreen,fonttitle=\bfseries,boxrule=0.7pt,arc=2mm,left=2.5mm,right=2.5mm,top=1mm,bottom=1mm,#1}
\newtcolorbox{reviewbox}[1][]{enhanced,breakable,colback=softorange,colframe=mainorange,coltitle=white,colbacktitle=mainorange,fonttitle=\bfseries,boxrule=0.7pt,arc=2mm,left=2.5mm,right=2.5mm,top=1mm,bottom=1mm,#1}
\newcommand{\chaptertitle}[1]{\vspace{0.5em}\begin{tcolorbox}[enhanced,colback=deepblue,colframe=deepblue,arc=2mm,boxrule=0pt]\begin{center}\Large\bfseries\color{white}#1\end{center}\end{tcolorbox}\vspace{0.5em}}
\newcommand{\sectitle}[1]{\vspace{0.5em}\noindent{\Large\bfseries\color{mainblue}#1}\par\vspace{0.25em}{\color{mainorange}\hrule height 1.2pt}\vspace{0.45em}}
\newcommand{\answerline}{\textcolor{answerred}{\rule{2.4cm}{0.7pt}}}

\begin{document}

\tableofcontents

\newpage

\chaptertitle{PPT 与章节测验综合背诵版（重排美化版）}

\begin{mainbox}

\textbf{使用方式：} 每章先背“课件重点”，再用“高频词”抓题眼，最后做“典型练习”。这一版按背诵节奏重排，不再堆满原始材料。

\end{mainbox}

\newpage

\section*{第一章土方工程}

\addcontentsline{toc}{section}{第一章土方工程}

\begin{mainbox}

\textbf{本章定位：} 共 35 道章节题。先抓概念边界，再背工艺顺序、质量控制和责任主体。

\end{mainbox}

\sectitle{一、PPT 重点压缩}

\begin{reviewbox}

\begin{enumerate}[label=\arabic*.]

\item 土方工程概述：一、土方工程概述

\item 土方工程概述：（一）土方工程的内容

\item 土方工程概述：包括一切土的挖掘、填筑、运输等过程以及排水降水、土壁支撑等准备工作和辅助工程。

\item 土方工程概述：（二）土的工程分类

\item 土方工程概述：水利水电工程施工组织设计规范

\item 土方工程的机械化施工：二、土方工程的机械化施工

\item 土方工程的机械化施工：（一）主要土方机械的特点与施工方法

\item 土方工程的机械化施工：、推土机

\end{enumerate}

\end{reviewbox}

\sectitle{二、题目高频词}

\begin{mainbox}

\textbf{看到这些词要立刻回忆：} 正确的有、正确的是、看答案、该题型暂不支持查、即进行下道工序时、地上十二层、地下一层、基坑挖好后不能立、当用人工挖土、的说法、错误的是、一个平面为矩形基

\end{mainbox}

\sectitle{三、典型练习题眼}

\begin{enumerate}[label=\arabic*.]

\item 第7题：土方机械的选择，一般应考虑的因素有。（答案：ACDE）

\item 第9题：以下关于内摩擦角的说法中，正确的是。（答案：ABCD）

\item 第14题：关于土方开挖的说法，正确的有。（答案：BCE）

\item 第17题：填方土料应符合设计要求，一般不能作为回填土的有。（答案：ABC）

\item 第19题：土方回填时关于基地处理以及土方填筑与压实的说法中，正确的有。（答案：ACD）

\item 第20题：关于土方回填的说法，正确的有。（答案：ACE）

\item 第25题：基坑开挖完毕后，必须参加现场验槽并签署意见的单位有。（答案：BCDE）

\item 第32题：关于土方填筑与压实的说法，正确的有。（答案：BD）

\end{enumerate}

\sectitle{四、背诵口令}

\begin{reviewbox}

\begin{itemize}

\item 概念题：先背定义，再背适用条件。

\item 工艺题：按施工顺序记，遇到“先后、分层、对称、及时”重点标记。

\item 质量安全题：按责任主体、检查方法、合格条件、整改要求四步背。

\end{itemize}

\end{reviewbox}

\newpage

\section*{第二章桩基础工程}

\addcontentsline{toc}{section}{第二章桩基础工程}

\begin{mainbox}

\textbf{本章定位：} 共 14 道章节题。先抓概念边界，再背工艺顺序、质量控制和责任主体。

\end{mainbox}

\sectitle{一、PPT 重点压缩}

\begin{reviewbox}

\begin{enumerate}[label=\arabic*.]

\item 地基处理技术与桩的类型：一、常用地基处理技术

\item 地基处理技术与桩的类型：（一）换填地基

\item 地基处理技术与桩的类型：换填地基

\item 地基处理技术与桩的类型：）换填厚度由

\item 地基处理技术与桩的类型：设计确定

\item 地基处理技术与桩的类型：，一般宜为

\item 地基处理技术与桩的类型：0.5

\item 钢筋混凝土预制桩：三、钢筋混凝土预制桩

\end{enumerate}

\end{reviewbox}

\sectitle{二、题目高频词}

\begin{mainbox}

\textbf{看到这些词要立刻回忆：} 正确的有、看答案、该题型暂不支持查、下列关于泥浆护壁、地下一层、总建筑面积、采用锤击沉桩法施、上部采用卡扣式接、下列关于静压桩终、下导管及桩身混凝、不同规格钢筋混凝、不妥之处

\end{mainbox}

\sectitle{三、典型练习题眼}

\begin{enumerate}[label=\arabic*.]

\item 第14题：某综合楼工程，地下三层，地上二十层，总建筑面积68000m³，地基基础设计等级为甲级，灌注桩筏板基础，现浇钢筋混凝土框架剪力墙结构，建设单位与施工单位按照《建设工程施工合同示范文本》GF-2013-0201签订了施工合同。基础设计桩径φ800mm，长度35～42m，混凝土强度等级C30，共计900根。施工单位编制的桩基施工方案中列明：采用泥浆护壁成孔、导管法水下灌注C30混凝土；灌注时桩顶混凝土面超过设计标高500mm；每根桩留置1组混凝土试件；成桩后按总桩数的20\%对桩身质量进行检验，监理工程师审查时认为方案存在错误，要求施工单位改正后重新上报。问题：指出桩基施工方案中的错误之处，并分别写出相应的正确做法。 该题型暂不支持查看答案。（答案：未标注）

\item 第10题：关于钢筋混凝土预制桩锤击沉桩顺序的说法，正确的有。（答案：ABC）

\item 第12题：某企业新建研发中心大楼工程，地下一层，地上十六层，总建筑面积 28000m2，基础为钢筋混凝土预制桩，二层以上为装配式混凝土结构，外墙装饰部分为玻璃幕墙，实行项目总承包管理。在静压预制桩施工时，桩基专业分包单位按照“先深后浅，先大后小，先长后短，先密后疏”的顺序进行，上部采用卡扣式接桩方法，接头高出地面0.8m。桩基施工后经检测，有1\%的II类桩。问题：桩基的沉桩顺序是否正确？卡扣式接桩高出地面0.8m是否妥当并说明理由？桩身的完整性有几类？写出II类桩的缺陷特征。 该题型暂不支持查看答案。（答案：未标注）

\item 第3题：下列关于泥浆护壁钻孔灌注桩施工要求的说法中，正确的有。（答案：BE）

\item 第4题：下列关于泥浆护壁灌注桩桩底注浆的说法中，正确的有。（答案：CE）

\item 第11题：关于泥浆护壁灌注桩的说法，正确的有。（答案：AC）

\item 第1题：预制桩基础施工中，桩起吊以及运输和施打时混凝土强度分别应达到设计强度的。（答案：D）

\item 第2题：下列关于静压桩终止沉桩标准的说法，错误的是。（答案：D）

\end{enumerate}

\sectitle{四、背诵口令}

\begin{reviewbox}

\begin{itemize}

\item 概念题：先背定义，再背适用条件。

\item 工艺题：按施工顺序记，遇到“先后、分层、对称、及时”重点标记。

\item 质量安全题：按责任主体、检查方法、合格条件、整改要求四步背。

\end{itemize}

\end{reviewbox}

\newpage

\section*{第三章基坑工程}

\addcontentsline{toc}{section}{第三章基坑工程}

\begin{mainbox}

\textbf{本章定位：} 共 23 道章节题。先抓概念边界，再背工艺顺序、质量控制和责任主体。

\end{mainbox}

\sectitle{一、PPT 重点压缩}

\begin{reviewbox}

\begin{enumerate}[label=\arabic*.]

\item 基坑支护施工技术：二、基坑支护施工技术

\item 基坑支护施工技术：施工前应具备的资料

\item 基坑支护施工技术：）岩土工程勘察报告、施工所需的设计文件；

\item 基坑支护施工技术：）施工影响范围内的建（构）筑物、地下管网和障碍物资料、施工组织设计、专项施工方案和施工监测方案。

\item 基坑支护施工技术：基坑工程专项施工方案的内容

\item 基坑支护施工技术：）支护结构、地下水控制、土方开挖和回填等施工技术参数；

\item 基坑支护施工技术：）基坑工程施工工艺流程；

\item 基坑监测：三、基坑监测

\end{enumerate}

\end{reviewbox}

\sectitle{二、题目高频词}

\begin{mainbox}

\textbf{看到这些词要立刻回忆：} 正确的有、以下关于地下连续、以下几种井点降水、的施工做法中、一般采用的降排水、下列关于内撑式排、下列关于悬臂式排、下水为目的的井是、不以降低基坑内地、不宜用于填土的是、井点回灌技术的主、以下关于土钉墙构

\end{mainbox}

\sectitle{三、典型练习题眼}

\begin{enumerate}[label=\arabic*.]

\item 第3题：下列关于内撑式排桩的说法中正确的有。（答案：ABD）

\item 第4题：地下连续墙施工之前应先设置的导墙的要求有。（答案：AB）

\item 第6题：以下关于地下连续墙混凝土导管法浇筑的说法中，正确的有。（答案：C）

\item 第7题：以下关于地下连续墙墙底注浆的说法中，正确的有。（答案：BE）

\item 第8题：以下关于土钉墙构造要求的说法中，正确的有。（答案：AE）

\item 第22题：应根据 等确定基坑工程监测项目、监测点布置、监测方法、监测频率和监测预警。（答案：BCD）

\item 第1题：下列关于悬臂式排桩的说法中错误的是。（答案：C）

\item 第5题：地下连续墙槽内泥浆面不应低于导墙面。（答案：A）

\end{enumerate}

\sectitle{四、背诵口令}

\begin{reviewbox}

\begin{itemize}

\item 概念题：先背定义，再背适用条件。

\item 工艺题：按施工顺序记，遇到“先后、分层、对称、及时”重点标记。

\item 质量安全题：按责任主体、检查方法、合格条件、整改要求四步背。

\end{itemize}

\end{reviewbox}

\newpage

\section*{第四章混凝土结构工程}

\addcontentsline{toc}{section}{第四章混凝土结构工程}

\begin{mainbox}

\textbf{本章定位：} 共 89 道章节题。先抓概念边界，再背工艺顺序、质量控制和责任主体。

\end{mainbox}

\sectitle{一、PPT 重点压缩}

\begin{reviewbox}

\begin{enumerate}[label=\arabic*.]

\item 混凝土工程：三、混凝土工程

\item 混凝土工程：（一）水泥的性能和应用

\item 混凝土工程：无机水硬性胶凝材料；

\item 混凝土工程：水硬性物质

\item 混凝土工程：硅酸盐水泥、铝酸盐水泥、硫铝酸盐水泥、氟铝酸盐水泥、磷酸盐水泥等；

\item 混凝土工程：用途性能

\item 大体积混凝土工程：四、大体积混凝土工程

\item 大体积混凝土工程：（一）施工组织设计

\end{enumerate}

\end{reviewbox}

\sectitle{二、题目高频词}

\begin{mainbox}

\textbf{看到这些词要立刻回忆：} 正确的是、某跨度、下列水泥品种中、下列钢材包含的化、下表、关于大体积混凝土、关于钢筋加工的说、其同条件养护试件、包括、可拆除该梁底模的、合要求的是、学元素中

\end{mainbox}

\sectitle{三、典型练习题眼}

\begin{enumerate}[label=\arabic*.]

\item 第15题：钢筋代换时应满足的构造要求有。（答案：CDE）

\item 第16题：钢筋代换时应满足的设计要求有。（答案：AB）

\item 第23题：模板工程设计的安全性原则是指模板要具有足够的。（答案：ACE）

\item 第32题：根据《通用硅酸盐水泥》GB175—2007 的规定，关于六大常用水泥凝结时间的说法，正确的是。（答案：B）

\item 第55题：下列关于砂浆、混凝土流动性说法中正确的有。（答案：BDE）

\item 第60题：高湿度环境中或长期处于水中的混凝土，应优先选用。（答案：BCDE）

\item 第65题：关于混凝土表面碳化的说法，正确的有。（答案：ABC）

\item 第75题：混凝土施工缝留置位置正确的有。（答案：ACDE）

\end{enumerate}

\sectitle{四、背诵口令}

\begin{reviewbox}

\begin{itemize}

\item 概念题：先背定义，再背适用条件。

\item 工艺题：按施工顺序记，遇到“先后、分层、对称、及时”重点标记。

\item 质量安全题：按责任主体、检查方法、合格条件、整改要求四步背。

\end{itemize}

\end{reviewbox}

\newpage

\section*{第五章预应力混凝土工程}

\addcontentsline{toc}{section}{第五章预应力混凝土工程}

\begin{mainbox}

\textbf{本章定位：} 共 7 道章节题。先抓概念边界，再背工艺顺序、质量控制和责任主体。

\end{mainbox}

\sectitle{一、PPT 重点压缩}

\begin{reviewbox}

\begin{enumerate}[label=\arabic*.]

\item 预应力混凝土概述：一、预应力混凝土概述

\item 预应力混凝土概述：（一）预应力混凝土工作原理

\item 预应力混凝土概述：（二）预应力混凝土的特点

\item 预应力混凝土概述：抗裂性好，刚度大；

\item 预应力混凝土概述：节省材料，减小自重；

\item 预应力混凝土概述：可以减小混凝土梁的竖向剪力和主拉应力；

\item 预应力混凝土概述：提高受压构件的稳定性；

\item 预应力施加方法：六、预应力施加方法

\end{enumerate}

\end{reviewbox}

\sectitle{二、题目高频词}

\begin{mainbox}

\textbf{看到这些词要立刻回忆：} 正确的是、下列预应力损失中、不应低于设计的混、为主、传递给混凝土、关于后张法预应力、关于预应力工程施、凝土立方体抗压强、力混凝土中的预应、包含的工序有、属于长期损失的是、工的说法

\end{mainbox}

\sectitle{三、典型练习题眼}

\begin{enumerate}[label=\arabic*.]

\item 第3题：关于预应力工程施工的说法，正确的是。（答案：D）

\item 第5题：关于后张法预应力施工的说法，正确的是。（答案：C）

\item 第7题：无粘结预应力施工包含的工序有。（答案：ACE）

\item 第1题：无粘结后张法预应力混凝土中的预应力，是靠 传递给混凝土。（答案：A）

\item 第2题：预应力筋放张，当设计无要求时，不应低于设计的混凝土立方体抗压强度标准值的。（答案：C）

\item 第4题：预应力筋的张拉以控制 为主。（答案：A）

\item 第6题：下列预应力损失中，属于长期损失的是。（答案：D）

\end{enumerate}

\sectitle{四、背诵口令}

\begin{reviewbox}

\begin{itemize}

\item 概念题：先背定义，再背适用条件。

\item 工艺题：按施工顺序记，遇到“先后、分层、对称、及时”重点标记。

\item 质量安全题：按责任主体、检查方法、合格条件、整改要求四步背。

\end{itemize}

\end{reviewbox}

\newpage

\section*{第六章砌筑工程}

\addcontentsline{toc}{section}{第六章砌筑工程}

\begin{mainbox}

\textbf{本章定位：} 共 31 道章节题。先抓概念边界，再背工艺顺序、质量控制和责任主体。

\end{mainbox}

\sectitle{一、PPT 重点压缩}

\begin{reviewbox}

\begin{enumerate}[label=\arabic*.]

\item 砌筑砂浆：一、砌筑砂浆

\item 砌筑砂浆：（一）原材料要求

\item 砌筑砂浆：原材料要求

\item 砌筑砂浆：）进场使用前应有出厂合格证和复试合格报告；

\item 砌筑砂浆：）强度等级——根据砂浆品种及强度等级要求选择；

\item 砖砌体工程：二、砖砌体工程

\item 砖砌体工程：（一）砌块种类

\item 砖砌体工程：粘土多孔砖

\end{enumerate}

\end{reviewbox}

\sectitle{二、题目高频词}

\begin{mainbox}

\textbf{看到这些词要立刻回忆：} 正确的是、正确的有、三一、下列做法中正确的、下列关于在砖墙上、下列关于混凝土小、下列关于砂浆搅拌、不应小于高度的、交接处应设置皮数、以下关于直槎处加、以下最适合用作砖、但直槎必须做成凸

\end{mainbox}

\sectitle{三、典型练习题眼}

\begin{enumerate}[label=\arabic*.]

\item 第8题：施工期间最高气温为25℃，砌筑用的普通水泥砂浆拌成后最迟必须在 内使用完毕。（答案：C）

\item 第14题：在砖砌体转角处、交接处应设置皮数杆，间距不宜大于。（答案：C）

\item 第21题：非抗震设防及抗震设防烈度为6度、7度地区的临时间断处，当不能留斜槎时，除转角处外，可留直槎，但直槎必须做成凸槎。以下关于直槎处加设拉结钢筋的说法中，正确的是。（答案：D）

\item 第23题：设有钢筋混凝土构造柱的抗震多层砖房，施工顺序正确的是。（答案：C）

\item 第9题：关于砌筑砂浆的说法，正确的有。（答案：ABC）

\item 第22题：对于设有钢筋混凝土构造柱的抗震多层砖房，下列做法中正确的是。（答案：BE）

\item 第25题：关于砌体结构工程施工的说法，正确的有。（答案：AD）

\item 第27题：关于普通混凝土小砌块施工说法，正确的是。（答案：BD）

\end{enumerate}

\sectitle{四、背诵口令}

\begin{reviewbox}

\begin{itemize}

\item 概念题：先背定义，再背适用条件。

\item 工艺题：按施工顺序记，遇到“先后、分层、对称、及时”重点标记。

\item 质量安全题：按责任主体、检查方法、合格条件、整改要求四步背。

\end{itemize}

\end{reviewbox}

\newpage

\section*{第七章脚手架工程}

\addcontentsline{toc}{section}{第七章脚手架工程}

\begin{mainbox}

\textbf{本章定位：} 共 10 道章节题。先抓概念边界，再背工艺顺序、质量控制和责任主体。

\end{mainbox}

\sectitle{一、PPT 重点压缩}

\begin{reviewbox}

\begin{enumerate}[label=\arabic*.]

\item 脚手架的种类：一、脚手架的种类

\item 脚手架的种类：它是为保证高处作业安全、顺利进行施工而搭设的工作平台或作业通道。在结构施工、装修施工和设备管道的安装施工中，都需要按照操作要求搭设脚手架。

\item 脚手架的种类：按其搭设位置分

\item 脚手架的种类：外脚手架和里脚手架；

\item 脚手架的种类：按其所用材料分

\item 脚手架的种类：木脚手架、竹脚手架与金属脚手架；

\item 脚手架工程的安全技术要求：三、脚手架工程的安全技术要求

\item 脚手架工程的安全技术要求：（一）脚手架的施工准备工作

\end{enumerate}

\end{reviewbox}

\sectitle{二、题目高频词}

\begin{mainbox}

\textbf{看到这些词要立刻回忆：} 其专项施工方案必、须进行专家论证的、下列分部分项工程、的有、需要进行专家论证、下列模板工程及支、下列脚手架工程中、以下关于脚手架搭、其专项施工方案不、卸工程中、取施工许可证或办、安全专项施工方案

\end{mainbox}

\sectitle{三、典型练习题眼}

\begin{enumerate}[label=\arabic*.]

\item 第2题：建设单位在申请领取施工许可证或办理安全监督手续时，应当提供。（答案：CE）

\item 第1题：以下关于脚手架搭设的说法中，正确的有。（答案：BC）

\item 第4题：下列模板工程及支撑系统工程中，其专项施工方案必须进行专家论证的有。（答案：AB）

\item 第5题：下列脚手架工程中，其专项施工方案必须进行专家论证的有。（答案：AE）

\item 第6题：起重吊装及安装拆卸工程中，其专项施工方案必须进行专家论证的有。（答案：BD）

\item 第7题：下列分部分项工程中，其专项施工方案必须进行专家论证的有。（答案：ABCD）

\item 第8题：下列分部分项工程的专项方案中，必须进行专家论证的有。（答案：ABDE）

\item 第9题：需要组织专家进行安全专项施工方案论证的是。（答案：B）

\end{enumerate}

\sectitle{四、背诵口令}

\begin{reviewbox}

\begin{itemize}

\item 概念题：先背定义，再背适用条件。

\item 工艺题：按施工顺序记，遇到“先后、分层、对称、及时”重点标记。

\item 质量安全题：按责任主体、检查方法、合格条件、整改要求四步背。

\end{itemize}

\end{reviewbox}

\newpage

\section*{第八章装饰工程（抹灰工程）}

\addcontentsline{toc}{section}{第八章装饰工程（抹灰工程）}

\begin{mainbox}

\textbf{本章定位：} 共 2 道章节题。先抓概念边界，再背工艺顺序、质量控制和责任主体。

\end{mainbox}

\sectitle{一、PPT 重点压缩}

\begin{reviewbox}

\begin{enumerate}[label=\arabic*.]

\item 抹灰工程施工：二、抹灰工程施工

\item 抹灰工程施工：（一）抹灰工程分类与材料的技术要求

\item 抹灰工程施工：）一般抹灰：水泥砂浆、水泥混合砂浆、聚合物水泥砂浆、粉刷石膏等抹灰。

\item 抹灰工程施工：）保温层薄抹灰：保温层外面聚合物砂浆薄抹灰。

\item 抹灰工程施工：）装饰抹灰：水刷石、斩假石、干粘石、假面砖等抹灰。

\item 抹灰工程施工：）清水砌体勾缝：清水砌体砂浆勾缝、原浆勾缝。

\item 拓展：装饰工程：章 装饰工程

\item 拓展：装饰工程：一、建筑装饰装修材料

\end{enumerate}

\end{reviewbox}

\sectitle{二、题目高频词}

\begin{mainbox}

\textbf{看到这些词要立刻回忆：} 关于抹灰工程的做、关于抹灰工程的施、工做法、正确的有、错误的有

\end{mainbox}

\sectitle{三、典型练习题眼}

\begin{enumerate}[label=\arabic*.]

\item 第1题：关于抹灰工程的做法，正确的有。（答案：BDE）

\item 第2题：关于抹灰工程的施工做法，错误的有。（答案：C）

\end{enumerate}

\sectitle{四、背诵口令}

\begin{reviewbox}

\begin{itemize}

\item 概念题：先背定义，再背适用条件。

\item 工艺题：按施工顺序记，遇到“先后、分层、对称、及时”重点标记。

\item 质量安全题：按责任主体、检查方法、合格条件、整改要求四步背。

\end{itemize}

\end{reviewbox}

\newpage

\section*{第九章防水工程}

\addcontentsline{toc}{section}{第九章防水工程}

\begin{mainbox}

\textbf{本章定位：} 共 7 道章节题。先抓概念边界，再背工艺顺序、质量控制和责任主体。

\end{mainbox}

\sectitle{一、PPT 重点压缩}

\begin{reviewbox}

\begin{enumerate}[label=\arabic*.]

\item 建筑防水材料的特性与应用：一、建筑防水材料的特性与应用

\item 建筑防水材料的特性与应用：建筑防水

\item 建筑防水材料的特性与应用：构造防水

\item 建筑防水材料的特性与应用：材料防水

\item 建筑防水材料的特性与应用：刚性防水

\item 建筑防水材料的特性与应用：柔性防水

\item 建筑防水材料的特性与应用：防水砂浆

\item 建筑防水材料的特性与应用：防水混凝土

\end{enumerate}

\end{reviewbox}

\sectitle{二、题目高频词}

\begin{mainbox}

\textbf{看到这些词要立刻回忆：} 正确的有、下列各厚度改性沥、严禁采用热熔法施、关于卷材防水层搭、关于屋面卷材防水、关于屋面涂膜防水、关于屋面防水工程、关于重要建筑屋面、层施工工艺的说法、屋面防水施工基本、工的是、接缝的做法

\end{mainbox}

\sectitle{三、典型练习题眼}

\begin{enumerate}[label=\arabic*.]

\item 第1题：关于重要建筑屋面防水等级和设防要求的说法，正确的有。（答案：A）

\item 第2题：屋面防水施工基本要求正确的是。（答案：BDE）

\item 第3题：关于屋面防水工程的做法正确的有。（答案：BE）

\item 第4题：关于屋面卷材防水施工要求的说法，正确的有。（答案：ABCE）

\item 第6题：关于卷材防水层搭接缝的做法，正确的有。（答案：ACE）

\item 第7题：关于屋面涂膜防水层施工工艺的说法，正确的是。（答案：B）

\item 第5题：下列各厚度改性沥青防水卷材中，严禁采用热熔法施工的是。（答案：A）

\end{enumerate}

\sectitle{四、背诵口令}

\begin{reviewbox}

\begin{itemize}

\item 概念题：先背定义，再背适用条件。

\item 工艺题：按施工顺序记，遇到“先后、分层、对称、及时”重点标记。

\item 质量安全题：按责任主体、检查方法、合格条件、整改要求四步背。

\end{itemize}

\end{reviewbox}

\newpage

\section*{第十章流水施工原理}

\addcontentsline{toc}{section}{第十章流水施工原理}

\begin{mainbox}

\textbf{本章定位：} 共 23 道章节题。先抓概念边界，再背工艺顺序、质量控制和责任主体。

\end{mainbox}

\sectitle{一、PPT 重点压缩}

\begin{reviewbox}

\begin{enumerate}[label=\arabic*.]

\item 有节奏流水施工：（七）流水施工基本组织方式

\item 有节奏流水施工：流水施工

\item 有节奏流水施工：非（无）节奏流水施工

\item 有节奏流水施工：有节奏

\item 有节奏流水施工：等节奏（固定节拍）流水施工

\item 有节奏流水施工：异节奏流水施工

\item 有节奏流水施工：异步距异节奏流水施工（一般的成倍节拍流水施工）

\item 有节奏流水施工：等步距异节奏流水施工（加快的成倍节拍流水施工）

\end{enumerate}

\end{reviewbox}

\sectitle{二、题目高频词}

\begin{mainbox}

\textbf{看到这些词要立刻回忆：} 个施工过程、分为、个施工段组织固定、施工时、节拍流水施工、下列流水施工参数、个施工段组织加快、各施工过程的流水、建设工程流水施工、某分项工程有、的成倍节拍流水施、的特点包括

\end{mainbox}

\sectitle{三、典型练习题眼}

\begin{enumerate}[label=\arabic*.]

\item 第1题：关于横道图进度计划特点的说法，正确的是。（答案：C）

\item 第2题：在有足够工作面和资源的前提下，施工工期最短的施工组织方式是。（答案：C）

\item 第10题：某工程有3个施工过程，分3个施工段组织固定节拍流水施工，流水节拍为2天。各施工过程之间存在2天的工艺间歇时间，则流水施工工期为 天。（答案：C）

\item 第11题：某分项工程有8个施工过程，分为3个施工段组织固定节拍流水施工，各施工过程的流水节拍均为4天，第三与第四施工过程之间工艺间歇为5天，该工程工期是 天。（答案：D）

\item 第12题：某工程由5个施工过程组成，分为3个施工段组织固定节拍流水施工，在不考虑提前插入时间的情况下，要求流水施工工期不超过44天，则流水节拍的最大值为 天。（答案：C）

\item 第13题：某分部工程有3个施工过程，分为4个施工段组织加快的成倍节拍流水施工，各施工过程流水节拍分别是6天、6天、9天，则该分部工程的流水施工工期是 天。（答案：B）

\item 第14题：某分项工程有 4 个施工过程，分为 3 个施工段组织加快的成倍节拍流水施工，各施工过程的流水节拍分别为 4d、 8d、 2d 和 4d，则应组织 个专业工作队。（答案：C）

\item 第18题：下列施工过程中，必须列入施工进度计划，并且大多作为主导的施工过程或关键工作的有。（答案：ABD）

\end{enumerate}

\sectitle{四、背诵口令}

\begin{reviewbox}

\begin{itemize}

\item 概念题：先背定义，再背适用条件。

\item 工艺题：按施工顺序记，遇到“先后、分层、对称、及时”重点标记。

\item 质量安全题：按责任主体、检查方法、合格条件、整改要求四步背。

\end{itemize}

\end{reviewbox}

\newpage

\section*{第十一章网络计划技术}

\addcontentsline{toc}{section}{第十一章网络计划技术}

\begin{mainbox}

\textbf{本章定位：} 共 23 道章节题。先抓概念边界，再背工艺顺序、质量控制和责任主体。

\end{mainbox}

\sectitle{二、题目高频词}

\begin{mainbox}

\textbf{看到这些词要立刻回忆：} 工作、见下图、计划如下图所示、则工作、某工程双代号网络、天和第、某双代号网络计划、某工程双代号时标、网络计划如下图所、下列网络图中的绘、下图所示、两项工作

\end{mainbox}

\sectitle{三、典型练习题眼}

\begin{enumerate}[label=\arabic*.]

\item 第3题：根据网络计划时间参数计算得到的工期称之为。（答案：B）

\item 第5题：某工作持续时间2天，有两项紧前工作和三项紧后工作，紧前工作的最早开始时间分别是第3天、第6天计算坐标系，对应的持续时间分别是5天，1天。紧后工作的最早开始时间分别是第15天、第17天、第19天，对应的总时差分别是3天、2天、0天。该工作的总时差是 天。（答案：A）

\item 第22题：某双代号网络计划如下图（图中粗实线为关键工作），若计划工期等于计算工期，则自由时差一定等于总时差且不为零的工作有（）。（见下图） 某双代号网络计划如下图（图中粗实线为关键工作），若计划工期等于计算工期，则自由时差一定等于总时差且不为零的工作有（）。（答案：A）

\begin{center}

\includegraphics[width=0.82\linewidth,height=0.24\textheight,keepaspectratio]{images/6708446672130686763.jpg}

\end{center}

\item 第2题：工程网络计划中，工作的最迟开始时间是指在不影响 的前提下，必须开始的最迟时刻。（答案：C）

\item 第4题：双代号网络计划中，某工作最早第3天开始，工作持续时间2天，有且仅有2个紧后工作，紧后工作最早开始时间分别是第5天和第6天，对应总时差是4天和2天。该工作的总时差和自由时差分别是。（答案：A）

\item 第10题：关于双代号时标网络计划的说法，正确的是。（答案：B）

\item 第11题：关于双代号早时标网络计划的说法，正确的是。（答案：D）

\item 第12题：关于双代号网络计划中线路的说法，正确的有。（答案：BC）

\end{enumerate}

\sectitle{四、背诵口令}

\begin{reviewbox}

\begin{itemize}

\item 概念题：先背定义，再背适用条件。

\item 工艺题：按施工顺序记，遇到“先后、分层、对称、及时”重点标记。

\item 质量安全题：按责任主体、检查方法、合格条件、整改要求四步背。

\end{itemize}

\end{reviewbox}

\newpage

\end{document}